\subsection{Calibration selection is unstable at the reference}
\label{sec:selection-stability}
In complete sales at 220/220, the Ratio and Descending candidates have minimum calibration exposure coverages 91.02768\% and 91.00063\%: a margin of only 0.02704 points. We resample 1,829 item clusters jointly across A/B and rerun all five threshold searches and both utility choices in each of 200 draws, holding fits, scores, cap and $T$ fixed. Exposure utility selects Ratio in 103 draws (51.5\%) and Descending in 97 (48.5\%); case utility selects Excess in 96\% and Mixed in 4\%. No draw has an empty menu. The signed Ratio-minus-Descending calibration margin has central 95\% resampling range [-0.479,0.281] points. Thus the family choice at this reference is unstable under calibration resampling. These frequencies are conditional sensitivity summaries, not probabilities of selecting a population-optimal rule.

The corresponding Bike day-cluster redesign selects Ratio/Mixed for exposure in 96\%/4\% of 200 draws. A separate 4,000-draw fixed-threshold diagnostic covers all four complete-sales head settings and Bike; its eligibility failures must not be read as full-procedure failures. Appendix~\ref{app:v7-stability} reports every candidate, signed margin, and frequency. Head interventions alone would not establish a small selection margin; this audit measures one directly.

\subsection{Complete-sales A design and separate-window risk screening}
\label{sec:complete-screen}
We apply the archived screen structure to all four complete-sales head settings, fixing reporting cap $r=.632683$ and designing thresholds on A alone at $.95r=.601049$. The 20-candidate family is four head settings times five rankings. Each fixed candidate is checked on B using an item-bootstrap upper quantile of signed excess at tail $.05/20$ (4,000 draws). Among passers, each utility chooses using A coverage alone; choices are frozen before rereading the already-consumed Later outcomes. This is retrospective risk evidence, not a new confirmation trial.

All 16 A-feasible candidates pass B; the four Error candidates are infeasible. Thus B removes no additional candidate, and the changed outcomes cannot be attributed to screening alone. Every cell selects Excess/Descending. On Later, all eight selected condition records have pointwise 95\% risk intervals below $r$; even their 20-candidate-adjusted upper summaries are at most .62445 (Table~\ref{tab:v7-screen-main}). These eight records contain six distinct policies, not independent replications. The exchange remains: 35.83--36.39 case points forgone for 12.77--14.29 exposure points gained. At 220/220, the stricter procedure costs the case choice 12.77 case/10.14 exposure points versus its original choice; the exposure choice loses 33.34/6.78 points. Appendix~\ref{app:v7-screen} retains every candidate, interval and cost. The bootstrap is approximate, conditional on existing fits, and does not remove shared calendar shocks or establish a population guarantee.
\begin{table}[t]\centering\small
\caption{Complete-sales A-design/B-screen procedure at fixed reporting cap .63268. All four cells select Excess for cases and Descending for exposure. Coverage differences (pp) and upper risk summaries concern the previously evaluated Later window. $U_{20}$ is the approximate one-sided bootstrap upper quantile with tail .05/20; it is not a finite-sample guarantee.}\label{tab:v7-screen-main}
\begin{tabular}{lrrrr}\toprule
Trees $e/w$ & $\Delta c$ & $\Delta d$ & Case $U_{20}$ & Exposure $U_{20}$\\\midrule
220/220 & -36.39 & +13.08 & 0.62295 & 0.62445 \\
660/220 & -35.83 & +14.29 & 0.62306 & 0.62445 \\
220/660 & -36.02 & +12.77 & 0.62305 & 0.62442 \\
660/660 & -36.29 & +13.53 & 0.62387 & 0.62442 \\
\bottomrule\end{tabular}\end{table}

